\begin{abstract}
\large \textbf{摘要.}
\normalsize
长时程智能体任务的后训练在计算效率与泛化能力之间存在张力。监督微调（SFT）虽然计算高效，却往往带来域外（OOD）性能退化；端到端强化学习（E2E RL）能够保留 OOD 能力，但由于需要进行多轮 on-policy rollout，计算成本很高。我们提出 \textbf{PivotRL}，一种直接作用于已有 SFT 轨迹的新框架，试图同时获得 SFT 的计算效率与 E2E RL 的 OOD 准确性。PivotRL 依赖两个核心机制：其一，在局部执行 on-policy rollout，并筛选出 \emph{pivot}，即那些采样动作结果方差较高、信息量更大的中间轮次；其二，对功能等价的动作给予奖励，而不是要求与 SFT 演示字符串严格一致。我们从理论上证明，这两点会带来更强的学习信号和更高的自然梯度范数，同时尽可能保持与训练任务无关动作上的策略概率排序不变。与在相同数据上训练的标准 SFT 相比，PivotRL 在四个智能体领域上的域内平均准确率高出 $+4.17\%$，在非智能体任务上的 OOD 准确率高出 $+10.04\%$。特别是在智能体编程任务上，PivotRL 仅用 E2E RL 四分之一的 rollout 轮次数，就达到了有竞争力的准确率。PivotRL 已被用于 NVIDIA 的 Nemotron-3-Super-120B-A12B，并成为其生产级智能体后训练的核心方案。
\end{abstract}
